\subsection*{Question 2.3 - Input transformation $u = L f_T$}

\paragraph{Step 1 - Building $L$.}
The new input stacks the total thrust $T$ on top of the torque vector $n$.
The total thrust is $T = f_1 + f_2 + f_3 + f_4 = \begin{bmatrix} 1 & 1 & 1 & 1 \end{bmatrix} f_T$, and the torque is $n = P f_T$ from Question 2.1.
Stacking the two,
\begin{equation*}
u = \begin{bmatrix} T \\ n \end{bmatrix} = L f_T, \qquad
L = \begin{bmatrix} 1 & 1 & 1 & 1 \\ 0 & r & 0 & -r \\ r & 0 & -r & 0 \\ c & -c & c & -c \end{bmatrix} \in \mathbb{R}^{4\times 4}.
\end{equation*}

\paragraph{Step 2 - Why this transformation is useful.}
The rigid-body model only depends on the rotor thrusts through two quantities.
The total thrust $T$ enters the translational dynamics, where it acts along $z_B$, and the torque $n$ enters the rotational dynamics.
In terms of the individual thrusts $f_i$ the problem is coupled: changing a single $f_i$ changes the total thrust and also one or two components of the torque.
In terms of $u$ the problem is decoupled into four virtual channels: one for the thrust (height) and one for each of the roll, pitch and yaw torques.
This has several advantages:
\begin{itemize}
\item The attitude and height controllers can be designed directly in terms of $T$ and $n$, each acting on its own channel, which is what the rest of the lab does.
\item The transformation is linear, constant and invertible, so it introduces no extra dynamics or nonlinearity, and the controllers can ignore the rotor layout and the constants $r$ and $c$.
\item The physical commands for the four rotors are only recovered at the very end, through $L^{-1}$. This stage is known as the control allocation or mixer of a flight controller.
\end{itemize}
The only caveat is that not every $u$ can be realized, since each $f_i$ must be nonnegative and is bounded by the motors.
So the mixer can saturate when a large torque is requested close to zero or to full thrust.

\paragraph{Step 3 - Inverting the transformation.}
We solve $L f_T = u$ for $f_T$ by combining pairs of equations.
The four equations are
\begin{equation*}
f_1 + f_2 + f_3 + f_4 = T, \quad
r(f_2 - f_4) = n_x, \quad
r(f_1 - f_3) = n_y, \quad
c(f_1 - f_2 + f_3 - f_4) = n_z .
\end{equation*}
Adding the first and last equations (the latter divided by $c$) isolates the sum of rotors 1 and 3, and subtracting them isolates the sum of rotors 2 and 4:
\begin{equation*}
f_1 + f_3 = \frac{T}{2} + \frac{n_z}{2c}, \qquad
f_2 + f_4 = \frac{T}{2} - \frac{n_z}{2c}.
\end{equation*}
The second and third equations give the differences $f_2 - f_4 = n_x / r$ and $f_1 - f_3 = n_y / r$.
Adding and subtracting each sum with the corresponding difference yields
\begin{equation*}
f_1 = \frac{T}{4} + \frac{n_y}{2r} + \frac{n_z}{4c}, \qquad
f_2 = \frac{T}{4} + \frac{n_x}{2r} - \frac{n_z}{4c},
\end{equation*}
\begin{equation*}
f_3 = \frac{T}{4} - \frac{n_y}{2r} + \frac{n_z}{4c}, \qquad
f_4 = \frac{T}{4} - \frac{n_x}{2r} - \frac{n_z}{4c}.
\end{equation*}
In matrix form, $f_T = L^{-1} u$ with
\begin{equation*}
L^{-1} = \begin{bmatrix}
\frac{1}{4} & 0 & \frac{1}{2r} & \frac{1}{4c} \\[4pt]
\frac{1}{4} & \frac{1}{2r} & 0 & -\frac{1}{4c} \\[4pt]
\frac{1}{4} & 0 & -\frac{1}{2r} & \frac{1}{4c} \\[4pt]
\frac{1}{4} & -\frac{1}{2r} & 0 & -\frac{1}{4c}
\end{bmatrix}.
\end{equation*}
The existence of this explicit solution also proves that $L$ is invertible whenever $r \neq 0$ and $c \neq 0$.
The formula has a simple reading: each rotor gets a quarter of the total thrust, plus or minus the amounts needed to produce the requested roll, pitch and yaw torques.
